\section{Wer mehr wissen will}
\subsection*{Wer mehr wissen will}

\blackframe

% Signing
% Wie weiß ich, dass die Mail von dem kommt, der er sagt?
%	EMail verfassen
%		-> Mail Signing
%		(unvertraute korrekte unterschrift)
%	Schlüsselverwaltung
%		-> Fingerprint prüfen
%		-> Key Signing

\begin{frame}{Verschlüsselte E-Mail an uns \dots}{\dots mit Fingerprints}
	\begin{block}{\url{info@michael-herbst.com}}
		\begin{center}
			\tt
			\larger
			3C0D 5B31 CDC8 F3B8 52AD\\
			00E6 3887 DB67 44A9 DF9F
		\end{center}
	\end{block}
	\begin{block}{\url{jonas@letopolis.de}}
		\begin{center}
			\tt
			\larger
			9393 3D08 35F9 BC8A 02BB \\
			5D41 138F BE3C 3B0F 5529
		\end{center}
	\end{block}
\end{frame}

\blackframe

\begin{frame}{Mobile Geräte}
	\begin{itemize}
		\item K9 Mail (\url{https://github.com/k9mail/k-9/})
		\item OpenKeychain (\url{http://www.openkeychain.org})
		\item Anleitung: 
			\begin{itemize}
				\item (F-Droid von \url{https://f-droid.org/} installieren)
				\item K-9 und OpenKeychain herunterladen
				\item Opnekeychain starten und konfigurieren
			\end{itemize}
	\end{itemize}
\end{frame}

\begin{frame}{Facebook und PGP}
	\begin{itemize}
		\item \url{https://www.facebook.com/notes/protect-the-graph/securing-email-communications-from-facebook/1611941762379302}
		\item Key auf Profil hinzufügen
		\item Statusmeldungen auf Wunsch verschlüsselt
		\item Noch nicht voll funktionsfähig
		\item Nachteil: Facebook kennt alle IDs / E-Mail Addressen
		\item Lösung: Separaten Key nur für Facebook
	\end{itemize}
	\note[item]{Nur eine PR Masche}
\end{frame}

\begin{frame}{Outlook und Windows Mail}
	\begin{itemize}
		\item Gpg4win beinhalted GpgOL
		\item \url{http://gpg4win.de/handbuecher/einsteiger.html}
		\item Plugin für Outlook 2003--2013
		\item Nicht sehr komfortabel
		\item Windows Live Mail nicht unterstützt
	\end{itemize}
\end{frame}

\begin{frame}{Links und weitere Veranstaltungen}
	\begin{itemize}
		\item \url{http://cryptoparty-hd.de} \\[10pt]
		\item \textbf{Smartphone-Messenger} am Mi 09.12., 18:30, INF 368, Raum 432
		\item \textbf{Anonymität} am Di 12.01., 18:30, Bergheimer Str. 58, Raum 01.034
	\end{itemize}
\end{frame}

\begin{frame}{Bilder und Lizenzen}
	\smaller
	\begin{itemize}
		\item Von \url{https://en.wikipedia.org/wiki/File:Asymmetric_cryptography_-_step_2.svg}
			\begin{center} \includegraphics[height=0.7cm]{Bob} \hspace{10pt} \includegraphics[height=0.7cm]{Alice} \end{center}

		\item Vom ``Tango Icon Theme'' --- \url{http://tango.freedesktop.org} 
			\begin{center} \includegraphics[height=0.7cm]{Document}\hspace{10pt}  \includegraphics[height=0.7cm]{mail_open}\hspace{10pt}  \includegraphics[height=0.7cm]{mail_closed} \end{center}
		%\item Von der HGS Mathcomp --- \url{http://mathcomp.org}
		%	\begin{center} \includegraphics[height=0.7cm]{logoHGS} \end{center}
		\item Alle anderen Bilder sind eigene Werke oder Derivate
		\item Letztere werden unter CC by-sa 4.0 veröffentlicht
	
	\vspace \fill
	\begin{center}
		\small
		\includegraphics[scale=0.5]{cc-by-sa}\\
		\scalebox{0.7}{This work is licensed under a Creative Commons}\\[-4pt]
		\scalebox{0.7}{Attribution-ShareAlike 4.0 International Licence.}
	\end{center}
	\end{itemize}
\end{frame}
